\section{Introduction}\label{sec:intro}
Let $-D<0$ be a fundamental discriminant, and write $\OO_D$ for the
ring of integers of $\mathbb Q(\sqrt{-D})$.
Gauss's class number one problem asks when $\OO_D$ has unique factorization,
or equivalently when its ideal class number $h(-D)$ is one.
Heegner~\cite{heegner}, Baker~\cite{baker} and Stark~\cite{stark} proved that
this occurs exactly for
\begingroup
\abovedisplayskip=14pt
\belowdisplayskip=14pt
\abovedisplayshortskip=14pt
\belowdisplayshortskip=14pt
\[
 D\in\{3,4,7,8,11,19,43,67,163\}.
\]
\endgroup
Here we obtain a weaker bound by the determinant method, without using
Heegner's methods or the known classification.

\begin{theorem}\label{thm:main}
Let $-D<0$ be a fundamental discriminant with $h(-D)=1$.  Then
$D\leq331$.
\end{theorem}

The question is how small an explicit bound for a hypothetical tenth
class-number-one discriminant can be obtained from such an argument.  This
is the setting of the challenge~\cite{tao2026}, which uses the class
number one problem to test recent determinant arguments concerning
Landau--Siegel zeros~\cite{green2026,taosketch}.  Our earlier paper established
$D\leq498$ by comparing an upper bound for the norm of an interpolation
determinant with the prime powers forced to divide that norm~\cite{chung498}.
The present paper keeps that determinant and its spectral size estimate.
We sharpen the rank criterion by retaining the exact degree and vanishing
estimates in the earlier proof.  For a fixed polynomial space and fixed
derivative orders, this permits fewer evaluation points and improves the
upper bound for the determinant.

\subsection{Overview of the proof}\label{subsec:overview}
Suppose that $h(-D)=1$ and $D>331$.  Every element of $\OO_D$ has the form
$(a+b\sqrt{-D})/2$ with $a,b\in\mathbb Z$, so every non-rational element
has norm at least $D/4$.  If a prime $p$ splits or ramifies, an ideal of
norm $p$ must be principal.  Its generator cannot be a rational integer,
whose norm would be a square.  Thus every prime $p<D/4$ is inert.  In
particular, $D$ is odd, since otherwise $2$ ramifies.  Nor can $D$ be
composite: its least prime divisor would ramify and be at most
$\sqrt D<D/4$.  Since $-D$ is fundamental, the resulting odd prime $D$ is
congruent to $3$ modulo $4$; it cannot be $7$ modulo $8$, since then $2$
splits.  Hence $D$ is prime with $D\equiv3\pmod8$.

The bound $D\leq498$ from~\cite{chung498} now leaves only
\[
  D\in\{347,379,419,443,467,491\}.
\]
We exclude these six values using the determinant estimates
of~\cite{chung498} and the sharper rank criterion proved below.

We form the matrix by evaluating monomials $z^rw^s$ with $r+Ks<R$,
and their normalized $z$-derivatives, at pairs $(z,\bar z)$ with
$z\in\OO_D$.  The parameter $K$ determines the shape of this triangular
set of exponents.  At an inert prime $p$, conjugation agrees with
Frobenius, so $z^p-\bar z$ is divisible by $p$.  Expressing the monomials
in terms of powers of $z^p-w$ therefore gives prime powers dividing every
maximal minor, with a loss depending on the selected derivative orders.
If one of these minors is nonzero, its norm is a positive
integer with this prescribed divisibility.  On the other hand, choosing
the evaluation points near the origin gives an Archimedean upper bound
for the same norm.  We obtain a contradiction when the forced divisor
exceeds that upper bound.  The arithmetic estimate also uses larger inert primes: class number one
lets us bound the loss from non-inert primes by counting norm values.
The upper bound keeps the derivative columns at each node together,
using the spectral estimates of~\cite{chung498} to account for their
correlations.

The rank criterion determines how few points can be used while ensuring
that a nonzero maximal minor exists.  If the polynomial space has dimension
$M$, and we take $J$ derivative orders at each of $n$ points, the necessary
dimension inequality is $nJ\geq M$.  The sufficient inequality stated
in~\cite{chung498} does not hold uniformly over the auxiliary-prime
intervals used here. Its proof, however, contains more precise information.

As in that proof, we rescale the polynomial rows and derivative columns
by known powers of an auxiliary
inert prime before reducing modulo the prime.  A relation among the reduced
rows would give a polynomial in two variables with prescribed local
vanishing; its second variable records the powers of $z^p-w$.  Applying a differential operator and taking a determinant of
coefficients then produces a nonzero polynomial in one variable.  We bound
its degree and the multiplicity of its zeros, retaining their dependence
on the degree of the original relation in its second variable.  These
more precise bounds show that the total multiplicity exceeds the degree,
even when the earlier sufficient inequality does not apply.  This proves
full row rank for the smaller node counts.

We still need an auxiliary inert prime for which this rank criterion
applies.  The required primes are larger than $D/4$, where the elementary
argument no longer guarantees inertness.  Our chosen intervals lie above
$D$, so their primes cannot ramify.  If a prime in one of these intervals
were split, class number one would represent it by the norm form
\[
 p=x^2+xy+\frac{D+1}{4}y^2.
\]
Counting the possible representations bounds the number of split primes
in the interval.  When $y$ is even, the
norm is congruent to $x^2$ modulo $2$; an odd prime norm therefore requires
$x$ odd, so the admissible values of $x$ are spaced by two.  We use this parity refinement
for the two intervals needed when $K=17$.  Since the count is smaller than
the number of primes in the interval, at least one prime there is inert.
We also place each interval above twice the radius of the disk containing
the nodes.  Distinct nodes then remain distinct after reduction at every
prime in the interval, as the rank criterion requires.

For $D=491,467,443,419$ we use $K=43$.  For $D=379$ and $D=347$, the
stronger rank criterion permits $K=17$. At fixed $R$, reducing $K$ allows more powers of the conjugate
variable, changing both the divisibility and the Gram-matrix estimate. Their combined
comparison is more favorable at these two discriminants.  For each chosen polynomial space, the sharper criterion brings the
number of points close to the minimum $M/J$. With these
choices, the arithmetic lower bound exceeds the Meixner--Schur upper bound
for each of the six candidates.  All six values are therefore excluded,
proving Theorem~\ref{thm:main}.

The rest of the paper is organized around these steps.  Section~\ref{sec:comparison}
recalls the determinant inequality and states the imported size estimate.
Section~\ref{sec:rank} proves the exact rank criterion.
Section~\ref{sec:primes} supplies the auxiliary inert primes and fixes
the three parameter families, while Section~\ref{sec:completion} gives
the six numerical comparisons.  The accompanying calculations
and source code are available at
\url{https://github.com/dualverse-ai/tenth_discriminant}.
